查询先选择 IVF 列表，再以主码筛选候选，最后精排保留项并返回 $k$ 个 ID 和分数。这三个决策使用不同分数；GT 则提供独立于搜索过程的第四种排序，用于评估：
\begin{align}
 s_i(q)&=\|\mu_i\|_2^2-2\langle q',\mu_i\rangle,
                   &&\text{centroid routing score},\nonumber\\
 \primarydist(q,x)&=\|q'-\widehat p(x)\|_2^2,
                   &&\text{primary score},\nonumber\\
 \refineddist(q,x)&=\|q'-\widehat p(x)-\widehat r(x)\|_2^2
                       +\|q'\|_2^2,
                   &&\text{refined ranking score},\nonumber\\
 \rawdist(q,x)&=\|q-x\|_2^2,
                   &&\text{raw-space distance}.
 \label{eq:5-ranking-objects}
\end{align}
其中 $\mu_i$ 是变换后的质心，$\widehat p,\widehat r$ 是解码后的主码与残差重建。精排分数在重建平方距离上增加一个仅依赖查询的常数，因此对同一查询保持该距离的排序（\cref{sec:5-6}）。原始向量召回率在第 7 章另行评估。

\par\medskip\noindent\begin{minipage}{\linewidth}
\captionsetup{hypcap=false}
\centering
\begin{tikzpicture}[font=\small\linespread{1}\selectfont,
 box/.style={draw=ThesisNavy!60,rounded corners=2pt,align=center,
 text width=3.7cm,minimum height=3.3cm,inner sep=4pt},
 arr/.style={-{Latex[length=2mm]},draw=ThesisNavy!80,thick}]
\node[box,fill=ThesisNavy!3] (route) at (0,0)
 {\textbf{Route: $s_i$}\\[4pt]Rotate the query\\Score centroids\\Select $n_{\rm probe}$ lists\\Visited candidates $U$};
\node[box,fill=ThesisNavy!3] (primary) at (5.45,0)
 {\textbf{Primary screen: $D_P$}\\[4pt]Build LUTs; score $U$\\[3pt]
  \colorbox{ThesisNavy!14}{\strut Retain primary top-$c$}};
\node[box,fill=ThesisNavy!14] (refine) at (10.9,0)
 {\textbf{Refine: $D_H$}\\[4pt]Read residual codes\\Score survivors\\Select final top-$k$\\Return IDs and scores};
\draw[arr] (route)--node[above,align=center,font=\footnotesize] {$n_c$\\candidates}(primary);
\draw[arr] (primary)--node[above,align=center,font=\footnotesize] {$c$\\survivors}(refine);
\node[draw=ThesisNavy!35,rounded corners=2pt,align=center,text width=14.5cm,inner sep=4pt] at (5.45,-2.4)
 {Evaluation: compare returned IDs with raw-vector ground truth};
\end{tikzpicture}

\captionof{figure}[查询流程]{查询流程.}
\label{fig:5-query-path}
\end{minipage}\par\medskip
路由定义访问空间 $U$，包含 $n_c$ 个候选出现实例。主码筛选最多保留 $c=\max\{k,\lceil\alpha k\rceil\}$ 项供残差评分。\cref{fig:5-query-path} 的阴影标出依赖该容量的工作：候选保留、残差评分与最终选择。在查询、索引及 probe 数固定时，路由集合和主码评分数量保持不变。下表对照各阶段读取的数据、工作与排序规则。

\par\medskip\noindent\begin{minipage}{\linewidth}
\captionsetup{hypcap=false}
\centering
\captionof{table}[查询阶段的选择约束]{查询阶段约束。$c_v=\min(c,n_c)$ 计数有效主码保留候选。精确性针对所说明的计算分数；路由与量化重建决定原始空间近似。}
\label{tab:5-selection-contracts}
\label{tab:5-query-stages}
\footnotesize\setlength{\tabcolsep}{3pt}
\begin{tabularx}{\linewidth}{@{}>{\raggedright\arraybackslash}p{1.7cm}>{\raggedright\arraybackslash}p{2.1cm}*{3}{>{\raggedright\arraybackslash}X}@{}}
\toprule
阶段 & 输入 $\to$ 输出 & 读取数据 & 主要工作 & 顺序／精确性 \\
\midrule
变换 & $q\to q'$ & 查询；变换 & 应用 $\Pi$ & 变换精度见 \cref{sec:7-precision} \\
Probe & $n_{\rm list}\to n_{\rm probe}$ 个列表 & FP32 质心；列表长度 & 评分；radix；排名计数 & $(s_i,i)$；精确 probe 选择 \\
LUT & $q'\to$ 表 & 主码本 & $MD_s2^{8/\glut}$ 项 & 实数算术分解；FP32 舍入 \\
主码，$\gscan=0$ & $n_c\to c_v$ & 打包主码；LUT & $O(n_cM)$ 评分；压缩 & 按距离精确选择；截断并列时保留候选可不同 \\
主码，$\gscan\geq1$ & $n_c\to c_v$ & 相同主码表示 & 评分；压缩；归并 & $(\primarydist,p)$；指定前提下精确 \\
精排 & $c_v\to c_v$ 个分数 & 残差码；校正 & $O(c_vd)$ & 计算的 $\refineddist$；FP32 舍入可随 lane 分组改变 \\
最终 & $c_v\to k$ & 精排分数；ID & $O(kc)$ 槽位扫描 & $(\refineddist,\mathrm{ID})$；有 $k$ 个有限有效项时精确 \\
\bottomrule
\end{tabularx}
\end{minipage}\par\medskip
评测搜索采用十六 lane 精排与自动扫描调度；它们不同于库的默认设置，配置见 \cref{sec:a-e1-protocol,tab:a-components}。后续各节依次说明查表、打包访问、分组选择、共享界、精排及 probe 选择（\cref{sec:5-2,sec:5-3,sec:5-4,sec:5-5,sec:5-6,sec:5-7}）。
